\documentclass[11pt]{article}

\usepackage[utf8]{inputenc}
\usepackage[T1]{fontenc}
\usepackage{lmodern}
\usepackage{microtype}
\usepackage{amsmath, amssymb}
\usepackage{geometry}
\geometry{a4paper, margin=1in}
\usepackage{xcolor}
\usepackage{hyperref}
\hypersetup{colorlinks=true, linkcolor=blue!50!black, citecolor=blue!50!black, urlcolor=blue!50!black}
\usepackage[backend=biber, style=numeric-comp, sorting=nyt]{biblatex}
\addbibresource{innovation_compression.bib}

\title{Chaos Makes Many, Compression Keeps Few:\\
Where Innovation Comes From, and the Human Part}
\author{Vasili Gavrilov\\
\small Independent researcher\\
\small ORCID: \href{https://orcid.org/0009-0007-9371-5994}{0009-0007-9371-5994}}
\date{\today}

\begin{document}
\maketitle

\begin{abstract}
Innovation admits a single description: chaos and recombination generate candidates without limit;
the few that survive are those that \emph{compress} a structure many share, which is also why they
spread; and a population sifts them in parallel. Assembled from established results, recombinant
growth, the bias--variance decomposition, and cultural evolution, the account splits the work into
\emph{variation}, now a cheap machine commodity, and the selecting \emph{prior}, the inductive bias
that decides what to keep. In one line: innovation is a high-conviction prior that resonates with
many.

The contribution is threefold. (i)~The \emph{assembly}, these results read as one process under a
single lens. (ii)~The \emph{grounding}: a learner gains nothing from data without a prior the data cannot supply
(No-Free-Lunch), and the classical limits of computation, the halting problem, G\"odel's
incompleteness, and Rice's theorem, one limit in many guises, show that no system certifies its own
foundations from within. A system can certainly \emph{change} its inductive bias, and meta-learning
does; what it cannot do is warrant the change without appeal to a further bias it also did not derive.
The argument offered here is therefore a \emph{regress}, formally motivated but not itself a theorem:
at every level the commitment is imported rather than produced, so the human role is a floor set by
the structure of justification, not a gap that scale will close. (iii)~A \emph{proposal}: a stable, human-controlled, key-addressed
ledger that holds the prior, its index being the bias itself, used both as a person's archive and as
the context supplied to a machine; a partial step toward the long-standing lifetime-archive problem. No individual theorem is offered as new.
\end{abstract}

\section{The picture in one breath}
Innovation looks like this. Something generates a huge number of candidate patterns. Almost all are
junk. A few are kept. The ones kept are the ones that are short to describe and that fit a structure
many people already share, so they catch on and spread. The making is cheap and can be mechanised.
The keeping needs a prior, a sense of which few are worth it, and that sense is the scarce part.

The rest of the essay walks through that sentence slowly: how the many are made (Section
\ref{sec:make}), why only a few survive (Section \ref{sec:few}), what ``survive'' means (Section
\ref{sec:compress}), who does the sifting (Section \ref{sec:crowd}), the human part (Section
\ref{sec:human}), why that part is a limit of principle and not a temporary gap (Section
\ref{sec:limits}), the tools that help (Section \ref{sec:tools}), and an old problem they touch
(Section \ref{sec:archive}). Section \ref{sec:new} says plainly what is new here, which is only the
assembly. The frame throughout is that intelligence is
the discovery of compressors, things that describe a lot with a little \cite{gavrilov_compressors}.

\section{Many are made}
\label{sec:make}
A simple rule, run over and over, produces patterns of
unlimited variety; this is what chaos and nonlinear dynamics do \cite{gavrilov_emergence}. And ideas
combine: an innovation is usually old pieces put together in a new way, and the number of possible
combinations grows exponentially with the number of pieces \cite{weitzman1998}. The supply of
candidates is therefore effectively endless. Generation is the easy part; selection is not.

\section{Few are kept}
\label{sec:few}
Out of that endless supply, only a small, countable set of forms come to matter, and they can often
be \emph{enumerated}.

\begin{itemize}
\item In biology, evolution searches an enormous space of body plans yet arrives again and again at
the same few solutions: the camera eye evolved independently dozens of times, powered flight in
insects, birds, and bats, the streamlined body in sharks, dolphins, and the extinct ichthyosaurs.
The space of possible forms is vast; the workable ones are few \cite{conwaymorris2003}.
\item In physics, one rule of minimization fixes a form across wildly different scales: raindrops,
bubbles, planets, and stars settle into the sphere, and packed cells tend to the hexagon, because few
shapes minimize surface or energy \cite{thompson1917}.
\item In catastrophe theory, the mathematics of how smooth systems can suddenly jump, Thom showed that for up to four
controlling factors there are exactly \textbf{seven} basic shapes of change, the seven elementary
catastrophes \cite{thom1975}; Arnold extended the same kind of classification to the caustics, the
bright curves light makes \cite{arnold1984}, and to the generic behaviours of differential equations
\cite{arnold_ode}.
\item In complexity science, Wolfram ran the simplest possible programs (cellular automata) by the thousands and found they
fall into just \textbf{four classes}; only a couple are genuinely rich, such as Rule~110 (which can
compute anything) and Rule~30 (used as a randomness source) \cite{wolfram2002}.
\item In statistical physics, wildly different materials near a critical point behave identically; they fall into
a small number of \emph{universality classes}, which is the renormalization-group fact already used
in the compressor account \cite{gavrilov_compressors}.
\end{itemize}

The lesson repeats across very different domains, and that it repeats, rather than that one mechanism
is at work, is the point: catastrophe theory is not social selection, nor is renormalization. What
they share is a \emph{shape}, a vast space collapsing to a few forms that a short description
captures, the infinite squeezed into a handful. That shared shape is compression, and the next
section argues it is also what decides which ideas are kept.

\section{What ``kept'' means: it compresses}
\label{sec:compress}
The kept ones are those \emph{short to say and wide in what they explain}, capturing a structure
shared across many cases. That is the definition of a good
compressor. In a formal theory of creativity, Schmidhuber makes this precise for what we call interesting or beautiful: a pattern is
interesting when noticing it lets you describe the world more briefly than before, a step of
compression progress \cite{schmidhuber2010}. A theorem that unifies scattered facts, a statistical model that captures a
process in a short formula and predicts its future to within a set error, a melody that
many ears find right, a format that addresses a record by a short key \cite{gavrilov_atree}: each
replaces a lot with a little. ``Resonates with many'' and ``compresses a shared structure'' are the
same statement seen from two sides; what cannot be compressed at all, an irreducible computation with
no shortcut shorter than running it \cite{wolfram2002}, is exactly what leaves no portable insight.

\section{Who keeps them: the crowd, in parallel}
\label{sec:crowd}
The sifting is not done by one judge. It is done by a whole population at once. A society is, in
effect, a machine with many processors running side by side (in computing terms, MIMD: many
instructions on many data at the same time). Candidate ideas spread through it like a contagion;
in cognitive anthropology, Sperber showed how ideas spread and settle on shared \emph{cultural attractors} \cite{sperber1996}, and, in the sociology of diffusion, Rogers mapped the familiar S-curve of how an
innovation spreads through a population \cite{rogers2003}. This selection across a population is the
old engine of cultural evolution, the inheritance of useful patterns without each person
rediscovering them \cite{boyd1985, dawkins1976}, and the same engine drives a language, across its
speakers, toward more learnable, compressible forms \cite{kirby2008}. A candidate is ``selected'' when it resonates,
which, by the previous section, means when it compresses something many already half-hold.

\section{The human part: the rare prior that points}
\label{sec:human}
As machines grow capable, the scarce part of innovation is no longer generation.

Split the work into two terms, the way statisticians split a model's error. One term is
\emph{variance}: how much you explore, how widely you search, how sensitive you are to the data. The
other is \emph{bias}: the assumptions you commit to before looking, your prior
\cite{geman1992}. Machines are very strong on variance. They generate and search through candidates
cheaply and in parallel, and a whole research programme, open-ended and novelty search, aims to make machines generate
variety on their own \cite{stanley2015}. Generation, the making of many, is becoming a machine commodity.

What machines do not get for free is the \emph{prior}: which of the endless candidates is worth
keeping here. That this cannot come from the data alone is a theorem of machine learning. Averaged over all possible
problems, no search method beats any other; all of a method's advantage comes from a prior matched to
the problem at hand \cite{wolpert1996}. The prior is information from outside the data. In machine learning this prior is made operational as
\emph{regularization}, a bias imposed before the data speaks: a penalty (ridge, \(L_2\); lasso,
\(L_1\)), or noise injected during training (jitter), which is equivalent to such a penalty. Supplying it,
the standing bet about what to keep and what to throw away, is the scarce contribution; the economics
of this, cheap machine prediction making human judgement more valuable, has been spelled out
directly \cite{agrawal2018}.

A good prior is expensive because it is paid for slowly: over evolutionary
time, at birth in the genome, across history in language and craft, and over a lifetime in skill and
taste. A cat's nervous system is a prior tuned tightly to its niche; a human's is looser and still
experimenting, which is why humans keep finding new things. The most striking cases are people whose
prior was unusually strong and idiosyncratic. Ramanujan, with little formal training, simply
\emph{saw} mathematical identities that were then checked and found true; Bach and Mozart heard
structure others did not, structure that is itself mathematical and well charted
\cite{voloshinov}.\footnote{The author trained in music (graduate of a music school).} Their leaps
were not blind guesses. This is the disagreement worth being
clear about: the standard story says creative variation is \emph{blind}, and the human only selects
afterwards \cite{campbell1960}. The cases above suggest the opposite: the rare prior makes the variation itself \emph{non-blind},
pointing the search before any selection. A defender of the blind view can answer that the variation
is merely internal, fast and unseen, and perhaps it is; but that does not touch the claim here, that
the prior is what aims the search, and the prior, not the search, is the scarce thing.

The definition follows. \textbf{Innovation is a high-conviction prior that resonates with
many.} A strong, committed, often contrarian prior, at the scale of whole fields a Kuhnian paradigm shift
\cite{kuhn1962}, is the generative bet. Resonance, the spread across the parallel crowd, is the verification that the bet
caught a real shared structure. One point of precision: ``high-conviction prior'' here means a
\emph{strong, narrow assumption} (a strong inductive bias), not the textbook ``high bias'' that means
plain error. The bet looks like error to contemporaries and turns out to fit a deeper structure. Nor
is such a prior a wager placed once and left alone. It is tightened by iteration from both sides,
held too narrowly until the world objects, loosened until it picks out nothing in particular, then
narrowed again. Overfitting and underfitting are the two cheap failures; the expensive thing between
them is, in the plainest available word, \emph{fitting}, and that bracketing is what the search
actually consists of \cite{gavrilov_compressors}.

Putting the halves together: machines supply cheap, abundant, increasingly open-ended
\emph{variation}; humans supply the rare, expensive, non-blind \emph{prior} that points it at the few
patterns that will compress and resonate. Working together, the two invent what neither does alone.

\section{Computational limits and the imported prior}
\label{sec:limits}
One might expect that a powerful enough machine, or a quantum one, will eventually remove every limit
and even supply its own prior. Two things have to be kept apart here, and this section keeps them
apart throughout: a set of theorems, which are settled, and an argument built on top of them, which is
not a theorem and is not offered as one. The theorems come first; they apply to any computer, human
brains included.

No procedure decides, for all programs, whether they halt: no general method can always tell, in
advance, whether a given program will eventually stop or run forever. This is the \emph{halting
problem}, and it is not an isolated curiosity. The same boundary appears, in guises that are
equivalent or inter-derivable, as G\"odel's incompleteness (no consistent system strong enough for
arithmetic proves every truth it can state, or its own consistency), as Rice's theorem (no algorithm
decides a non-trivial property of what a program computes), as the word problem for string-rewriting
(semi-Thue) systems, across the interchangeable models of computation (Church's \(\lambda\)-calculus,
the recursive functions, Turing machines), and in Wolfram's simplest programs once one of them is
universal. These are theorems of computability theory, not gaps that faster hardware will close; one standard text works
through them and their equivalences \cite{davis_sigal_weyuker, wolfram2002}. Because the same limit
recurs under every guise, it belongs to computation itself, not to one model, and nothing internal
escapes it by being faster or cleverer. Quantum
computing does not change this: a quantum computer can be simulated by an ordinary one, so it computes
no function a classical machine cannot, and it is not known to solve NP-complete problems
efficiently. It speeds up a few tasks; it does not remove the limits.

Two cautions, and then the limit of the argument itself. First, these limits do not show that human
minds exceed machines; the known attempt to argue that from G\"odel does not hold, and we do not use
it. The limits apply to people and machines alike. Second, what they do show is narrower: a system
does not \emph{certify} its own foundations from within. A formal system cannot prove its own
consistency, so the warrant for using it comes from outside; and a learner has no advantage without a
prior, which the data cannot supply (No-Free-Lunch).

The third caution runs against the argument of this section, and it is the important one. Nothing
above forbids a system from \emph{changing} its own inductive bias, and systems do exactly that:
meta-learning adjusts a learner's bias from experience across tasks, hierarchical models place priors
over priors and update them, learned optimizers acquire their own update rules, and architecture
search selects structural commitments a researcher would otherwise hand-write. Anyone claiming these
are impossible is refuted by practice, and the present argument does not claim it. What none of them
does is \emph{escape} the problem, because each shifts the commitment up a level rather than removing
it: the meta-learner's own bias, the hyperprior, the space the architecture search ranges over, none
of these is derived from the data either, and each was chosen by someone. What survives is therefore a
regress rather than a corollary of computability theory. At every level the commitment is imported,
and there is no level at which it is produced. The theorems above make the formal cousin of this exact
in the case of proof, where the impossibility of self-certification is settled; the extension to
inductive bias is an argument from the same shape, and is offered as one. It stands close to Hume's
problem of induction and to Goodman's new riddle, and it inherits their standing rather than improving
on it.

One further concession is owed on the No-Free-Lunch side. The theorem averages uniformly over all
problems, and a uniform distribution over target functions is itself a prior that nobody holds; real
problem distributions are highly structured, which is why practical learners work at all. No-Free-Lunch
therefore does not show that a good prior is hard to come by, only that it cannot be extracted from the
data alone and has to come from somewhere else. That is all this argument needs, and it is worth
claiming no more than that. \emph{Prior} here is meant broadly, not in the Bayesian sense of a probability distribution but as
whatever a system needs and cannot generate from within: in physics the initial and boundary
conditions, which the laws of motion never fix; in logic the axioms; in learning the inductive bias;
in an AI system the context that frames its task. These are one move, not four; only the narrow
Bayesian reading would split them. The dynamics may be universal, but what picks out the actual case
is supplied from outside the computation, not produced by more of it. The same data admits different
readings under different priors, and nothing within the data settles which is right.

This grounds Section~\ref{sec:human}, with the status just stated: the prior is scarce not only for
economic reasons but because nothing internal to a computation warrants it. Whoever supplies it, so
far mostly humans, supplies what no amount of compute produces for itself, though compute can and
does move the question up a level.

\section{The tools that help}
\label{sec:tools}
A prior is only worth having if it can be kept and reused. It is stored in three places: in genes, in
culture, and in external artifacts such as notes, code, and archives. External storage lets a prior
\emph{compose}, build on itself across people and years, which is why human knowledge accumulates and
a single animal's does not \cite{gavrilov_compressors}. Tools that pack accumulated knowledge into a
compact, recallable form therefore extend the prior. The author's AIS is one such tool \cite{gavrilov_ais}: a personal,
plain-text store in which a short key addresses a body of content, so that recalling by a short key
is itself a small compression \cite{gavrilov_atree}.

For this to work the store must be \emph{stable}: a volatile memory keeps no prior. Today's AI memory layers (per-session instruction
and skill files, vector stores) are largely volatile and ad hoc, which is why the author has called
for a stable, human-controlled ledger for an AI's working knowledge \cite{gavrilov_context_renorm}, in an open call to the model vendors \cite{gavrilov_context_call}, and offers a second version of that protocol built on AIS: a collection of
plain-text files, each reached by a short key that points to the data. The key does the compressing, a short address for a lot of content; the files
themselves stay plain, so no brittle fixed format is locked in and the way they are read can change
while the files endure. The content itself can sit anywhere, even on the open web, held by reference
such as a URL, so one store indexes a person's files and the web alike, and serves both that person's
notes and an AI's working knowledge. What is personal and scarce is not the content but the index
over it: the person's own associations, a living map of concepts grown to fit present understanding
rather than adopted ready-made, which is the bias of earlier sections made concrete. The content is
never copied: the index points to it by reference, even into shared, read-only stores, leaving the
data in place. And since the index is one person's prior, another's differs; the tool keeps each
person's own ordering of the world rather than collapsing many priors onto one ontology, and by
No-Free-Lunch that diversity is a resource, not a flaw to correct.

Its advantages over today's chat-managed memory are stability and human control. When a machine
compresses its own running context, it does so with a prior, but not the keeper's: the one its
designers and its training objective installed, held fixed and applied to every ledger alike. Asked
to recompress a human's ledger it keeps what fits that objective and drops what the keeper marked as
essential, unable to tell, from its own context, which items are load-bearing rather than restatable.
When a human decides what the ledger keeps and what it drops, that decision comes from outside the system,
directed by the prior, the same prior no computation derives for itself (Section~\ref{sec:limits}),
and is made to fit the current understanding. The store can be strictly
append-only, like a lifetime archive, or curated by hand; either way the compaction is human-directed,
not a rewrite the model performs on itself mid-conversation. Against a flat instruction file, a record
is also reached directly by its key, rather than carried whole in context.

\section{An old problem: the lifetime archive}
\label{sec:archive}
The aim is older than the tools. Recording and retrieving what one accumulates over a lifetime, a
personal archive, is a long-standing grand challenge: Bush's Memex \cite{bush1945}, and later one of
Gray's dozen long-term goals for computing, to record what a person sees and hears and return any
item on request \cite{gray2003}, and named again as the UK research community's ``Memories for Life'' challenge \cite{ohara2006}; to this day the challenge stands. The ledger here does not solve it either, and it leaves aside the
separate, hard problem of keeping file formats readable across decades \cite{rothenberg1995}; it offers one practical,
partial step, a stable, human-controlled, format-compatible store whose index is the person's own
associations. The author's archive tools were built in that direction, motivated by those goals.

This rests on tested practice rather than a new idea: the author has kept and used such plain-text
context indexes since the early 2000s, organizing a personal archive that has outlived several disk
failures without data loss, kept by plain redundant copies in sync. None of these tools manufacture the
prior; they lower the cost of keeping and reusing the scarce
thing.

\section{What is new here, and what is not}
\label{sec:new}
Honesty about contribution. Every block above is established and named: combinatorial and recombinant
growth \cite{weitzman1998}; the collapse to a few canonical forms
\cite{thom1975, arnold1984, wolfram2002, conwaymorris2003, thompson1917}; compression as the mark of interest
\cite{schmidhuber2010}; culture as parallel contagion with attractors \cite{sperber1996, rogers2003};
variation and selection \cite{campbell1960, boyd1985}; the bias--variance split \cite{geman1992}; the
limits of computation, Church--Turing, G\"odel, Rice \cite{davis_sigal_weyuker}, and the No-Free-Lunch
limit \cite{wolpert1996}; cheap prediction and scarce judgement \cite{agrawal2018}. The
two neighbours closest to this essay, and most important to credit, are \textbf{Sperber} (contagion
and cultural attractors) and \textbf{Schmidhuber} (interest as compression). The contribution claimed
here is the \emph{assembly}: viewing all of these as one process under a single lens, generation by
chaos, collapse and selection by compression, sifting across a parallel social machine, with the
human supplying the rare non-blind prior. The one move the author presses as his own, rather than
borrows, is the grounding of that prior in the limits of computation (Section~\ref{sec:limits}),
offered with the status that section states: not as a corollary of the theorems, but as a regress they
motivate. A system can move its inductive bias up a level, and meta-learning does; at no level is the
commitment produced rather than imported, and that is the floor under the human's role here. Alongside the assembly,
the essay makes one concrete proposal (Section~\ref{sec:tools}): a stable, human-controlled,
plain-text, key-addressed ledger, its content held by reference and its index the person's own map of
concepts, grown rather than adopted ready-made. The data structure is not new, associative and graph
memories exist; what is proposed is the specific human-curated, plain-text, reference-indexed form,
and the point that the index is the bias; the proposal is offered as a partial step toward a
long-recognized grand challenge, the lifetime personal archive (Section~\ref{sec:archive}). No
individual theorem is offered as new.

\section*{Closing}
A last observation, exact rather than decorative. To write this essay is to take many separate, named
results and squeeze them into one short structure. That act, replacing a lot with a little that fits,
is the very thing the essay says innovation is. The map is an instance of the thing it maps.

\printbibliography

\end{document}
